% !TeX root = ../教材题库_全部条目_转写工作稿.tex
% Source: XJ-B1-C06-HW-P254-Y0444
% 第6章 / 6.4 用样本估计总体 / 教材 p.254 / PDF 第261页
% OCR未校对：公式与文字均可能有误；未完成数学TeX转写；禁止直接发给学生
% 原题及相邻原解截图（核对用）：../题目原图/XJ-B1-C06-HW-P254-Y0444-01.webp
% 本文件只在转写完成后作为题干使用；OCR 条目尚不是可靠题干。
\begin{bankitem}{习题6.4 第1题}{第6章 / 6.4 用样本估计总体 / 教材 p.254 / PDF 第261页}{XJ-B1-C06-HW-P254-Y0444}{OCR未校对：公式与文字均可能有误；未完成数学TeX转写；禁止直接发给学生}
\begin{ocrdraft}
[教材p.254 / PDF第261页]
\par
1.唐代是我国封建社会的兴盛时期．从现存的唐代墓志
\par
中，我们了解到95个唐代妇女的初婚年龄，如下表所示：
\par
初婚年龄/岁13
\par
14
\par
17
\par
18
\par
19
\par
20
\par
10
\par
9
\par
14
\par
19
\par
11
\par
7
\par
初婚年龄/岁21
\par
22
\par
23
\par
24
\par
25
\par
29
\par
33
\par
-00000000
\par
（甘肃安西榆林窟壁
\par
（1）计算上述样本的平均初婚年龄，据此估计唐代妇女
\par
画，张大千临摹）
\par
的平均初婚年龄．
\par
（2）试比较唐代妇女平均初婚年龄与我国当代妇女法定初婚年龄的大小．
\par
（3）通过互联网调查世界各国妇女初婚年龄，妇女初婚年龄是否随着时代的发
\par
展而逐渐增大？请尝试解释说明．
\end{ocrdraft}
\end{bankitem}
